\documentclass[10pt,a4paper]{article}
\usepackage{fontspec,amsmath,amssymb,graphicx}
\usepackage[margin=18mm]{geometry}
\setmainfont{Times New Roman}
\setlength{\parindent}{0pt}
\setlength{\parskip}{5pt}
\allowdisplaybreaks
\sloppy
\emergencystretch=2em
\begin{document}
\section*{VD\_3\_1\_2\_01 — Lời giải bổ sung}
\subsection*{Kiến thức và trình tự giải}
\textbf{Bước 1 — Nhận dạng dữ liệu và tốc độ.} Viết $u_{tt}-a^2u_{xx}=F$, $u(x,0)=f(x)$, $u_t(x,0)=g(x)$; tốc độ là $a>0$, không phải $a^2$. Với dữ liệu trơn, cần $f\in C^2$, $g\in C^1$. Nếu dữ liệu có góc hoặc bước nhảy, công thức vẫn mô tả sóng từng miền nhưng phải xác định nơi mất tính trơn.

\textbf{Bước 2 — Tìm hai sóng chạy.} Đặt $\xi=x-at$, $\eta=x+at$. Quy tắc dây chuyền cho $u_{tt}-a^2u_{xx}=-4a^2u_{\xi\eta}$, nên khi $F=0$, $u=A(\xi)+B(\eta)$. Tại $t=0$, $A+B=f$, $a(B'-A')=g$. Chọn một nguyên hàm $H'=g$ rồi lấy $A=f/2-H/(2a)$, $B=f/2+H/(2a)$; hằng số của $H$ sẽ triệt tiêu.


\[
u_0(x,t)=\frac{f(x-at)+f(x+at)}2+\frac{H(x+at)-H(x-at)}{2a}.
\]


\textbf{Bước 3 — Tính tích phân theo dữ liệu thật.} Tìm nguyên hàm trên từng khoảng, nối liên tục nếu $g$ có bước nhảy; thay hai cận $x\pm at$. Dùng công thức cộng lượng giác trước khi kết luận chu kỳ. Với dữ liệu chỉ thị, tích phân bằng độ dài giao của $[x-at,x+at]$ với đoạn hỗ trợ, nhân hệ số của $g$.

\textbf{Bước 4 — Nếu có nguồn, thêm Duhamel.} Nghiệm riêng có dữ liệu đầu bằng không là


\[
u_F(x,t)=\frac1{2a}\int_0^t\int_{x-a(t-s)}^{x+a(t-s)}F(y,s)\,dy\,ds.
\]


Tính tích phân trong trước. Nếu $F$ là đa thức hoặc một mode lượng giác, có thể đoán $u_F$ rồi giải ODE theo $t$, nhưng vẫn phải buộc $u_F(x,0)=(u_F)_t(x,0)=0$. Nghiệm cuối là $u_0+u_F$.

\textbf{Bước 5 — Kiểm tra và đọc ý nghĩa.} Đạo hàm công thức: $u_x=(f'(\xi)+f'(\eta))/2+(g(\eta)-g(\xi))/(2a)$, $u_t=a(f'(\eta)-f'(\xi))/2+(g(\eta)+g(\xi))/2$. Tại $t=0$ thu lại $f,g$; đạo hàm tiếp cho PDE. Các điểm gãy $s=b$ phát sinh đường $x\pm at=b$. Chu kỳ phải đúng cho mọi $x,t$, không chỉ sau khi một mặt sóng đi qua một điểm cố định.
\subsection*{Lời giải áp dụng cho bài}
\textbf{Bước 1. Tách li độ và vận tốc.} Với $a=1$ ở (i), (iii), (iv), và $a=2$ ở (ii), dùng
\[u=\tfrac12[f(x-at)+f(x+at)]+\frac{H(x+at)-H(x-at)}{2a},\qquad H'=g.\]
\textbf{Bước 2. Tính từng trường hợp.}

(i) $g=0$, nên $H=0$. Công thức $\sin(x-t)+\sin(x+t)=2\sin x\cos t$ cho $u=\sin x\cos t$.

(ii) Chỉ đổi độ dịch thành $2t$: $u=\sin x\cos2t$. Hai chu kỳ cơ bản lần lượt là $2\pi$ và $\pi$.

(iii) $f=0$, $g(s)=s|s|$. Trên $s>0$, nguyên hàm là $s^3/3$; trên $s<0$, là $-s^3/3$. Nối tại $0$ được $H(s)=|s|^3/3$, do đó
\[u=\frac{|x+t|^3-|x-t|^3}{6}.\]
Hệ số $1/2$ ở nguyên hàm trong bản scan không đúng: $(|s|^3/3)'=s|s|$. Hàm $g$ thuộc $C^1$, nên nghiệm thuộc $C^2$, kể cả trên $x\pm t=0$; đạo hàm cấp ba có thể nhảy.

(iv) $g(s)=|s|$ có nguyên hàm $H=s|s|/2$, nên
\[u=\frac{(x+t)|x+t|-(x-t)|x-t|}{4}.\]
Vì $H\in C^1$ nhưng không $C^2$ tại $0$, nghiệm $C^1$ và không $C^2$ trên hai đặc trưng $x=\pm t$.

\textbf{Bước 3. Kiểm tra.} Tại $t=0$, hai hiệu nguyên hàm bằng không; đạo hàm theo $t$ cho $g(x)$. Mỗi hạng chỉ phụ thuộc $x\pm at$ nên thỏa PDE nơi đạo hàm tồn tại. Hai nghiệm (iii), (iv) tăng theo $t$ ở những điểm thích hợp, nên không tuần hoàn.\end{document}